\documentclass[11pt]{article}
\usepackage[margin=0.92in]{geometry}
\usepackage{amsmath,amssymb,amsthm,mathtools}
\usepackage{enumitem}
\usepackage{graphicx}
\usepackage{xcolor}
\usepackage{microtype}
\usepackage[hidelinks]{hyperref}

\newtheorem{theorem}{Theorem}
\newtheorem{lemma}[theorem]{Lemma}
\newtheorem{corollary}[theorem]{Corollary}
\newcommand{\D}{\mathbb D}
\newcommand{\T}{\mathbb T}
\newcommand{\C}{\mathbb C}
\newcommand{\N}{\mathbb N}
\newcommand{\Hp}{H^p}
\newcommand{\Hq}{H^q}
\newcommand{\Lin}{\mathcal L}
\newcommand{\Orb}{\operatorname{Orb}}
\newcommand{\supp}{\operatorname{supp}}

\title{A sharp Hilbert-core obstruction at \(p>2\)\\
for Shkarin's coanalytic Toeplitz question}
\author{Autonomous mathematical investigation, lane 7}
\date{17 August 2026}

\begin{document}
\maketitle

\begin{abstract}
Question 5.4 of Shkarin asks whether an expanding coanalytic Toeplitz operator
can be weakly hypercyclic on $H^p$ for some $p>2$.  We do not settle the
existence question.  We prove, however, that the Hilbertian mechanism behind
the positive $H^2$ theorem cannot work above the endpoint under \emph{any}
choice of invariant Hilbert core.  Precisely, if $T$ is the coanalytic
Toeplitz operator induced by a nonconstant $g\in H^\infty$ with
$g(\D)\cap\D=\varnothing$, $U$ is a Hilbert-space isometry without nonzero
finite-dimensional invariant subspaces, and a bounded map $J$ satisfies
$TJ=JU$, then $J=0$.  The proof first identifies every bounded forward orbit
with a boundary functional in $L^p$ on the contact set.  Wold decomposition
then reduces the shift part to an impossible $H^2\to L^p$ boundary embedding
and the unitary part to an impossible non-atomic $L^2\to L^p$ embedding.
The exponent contradiction is exactly $p/2>1$.
\end{abstract}

\section{Source question and status}

Let $M_g$ denote multiplication by $g$ on a Hardy space.  In the Hilbert
case its adjoint is the coanalytic Toeplitz operator $T_g^*$.  Shkarin proves
weak hypercyclicity on $H^2$ when $g(\D)\cap\D=\varnothing$ and the boundary
contact set $\{|g|=1\}$ has positive measure.  After observing that the proof
extends by comparison of topologies to $1\le p\le2$, the paper asks:

\begin{quote}
\textbf{Question 5.4.} Is there any $g\in H^\infty(\D)$ with
$g(\D)\cap\D=\varnothing$ such that $T_g^*$ is weakly hypercyclic as an
operator on $H^p(\D)$ for some $p>2$?
\end{quote}

\begin{center}
\includegraphics[width=0.94\textwidth]{figures/open_question_crop.png}
\end{center}

Fix $2<p<\infty$ and let $q=p/(p-1)\in(1,2)$.  Under standard Hardy
duality, the coanalytic Toeplitz operator on $H^p$ is conjugate-linearly
isometrically similar to the Banach adjoint
\[
 T=M_g'\in\Lin((H^q)^*).
\]
This avoids using Hilbert-adjoint notation outside $p=2$.  Conjugate-linear
similarity preserves weak density; when an intertwining Hilbert space occurs,
one passes to its conjugate complex structure.

\begin{theorem}[No isometric Hilbert core above the endpoint]
\label{thm:main}
Let $2<p<\infty$, $q=p/(p-1)$, and let $g\in H^\infty(\D)$ be nonconstant
with $g(\D)\cap\D=\varnothing$.  Put $X=(H^q)^*$ and $T=M_g'$ on $X$.
Let $U$ be an isometry on a Hilbert space $\mathcal H$ having no nonzero
finite-dimensional invariant subspace.  If $J:\mathcal H\to X$ is bounded
and
\[
 TJ=JU,
\]
then $J=0$.
\end{theorem}

\begin{corollary}[Obstruction to the source criterion]
\label{cor:criterion}
There is no nonzero $T$-invariant subspace $F\subseteq X$ carrying an inner
product for which
\begin{enumerate}[label=\textup{(\roman*)},nosep]
\item the Hilbert norm is stronger than the norm inherited from $X$, and
\item $T|_F$ extends to an isometry without nonzero finite-dimensional
invariant subspaces on the Hilbert completion of $F$.
\end{enumerate}
In particular, conditions \textup{(B1)--(B2)} of Theorem 1.2 in the source
paper can never hold for this operator on $H^p$ when $p>2$.
\end{corollary}

\section{Idea of the proof}

Write $h=1/g$.  It is an analytic self-map of the disc.  If $x\in X$ has
a bounded forward $T$-orbit, then
\[
 x(f)=(T^nx)(h^nf).
\]
As $n\to\infty$, the $H^q$ norm of $h^nf$ loses all boundary mass except on
$A=\{|h|=1\}$.  Thus $x$ is represented by an $L^p(A)$ boundary density.
Every vector in the range of an intertwiner $J$ has a bounded orbit, because
$U$ is an isometry.  Hence $J$ factors through a bounded map
$K:\mathcal H\to L^p(A)$ intertwining $U$ with multiplication by the
unimodular function $g|_A$.

The Wold decomposition leaves two cases.  A wandering vector in the shift
part has an orthogonal orbit, so $K$ would turn the $H^2$ norm of analytic
polynomials into an $L^p$ norm against a nonzero boundary measure.  Hardy peak
functions force that measure to satisfy $\nu(I)\lesssim |I|^{p/2}$ on every
arc; a partition into equal arcs gives $\nu=0$.  On the unitary part, spectral
calculus would yield $\|b\|_{L^p(\nu)}\lesssim\|b\|_{L^2(\mu)}$ with $\mu$
non-atomic.  Partitioning $\mu$ into equal-mass pieces again gives
$\nu(\T)\lesssim N^{1-p/2}\to0$.  Both Wold components vanish.

\section{Two concentration lemmas}

Normalized analytic peak functions give the endpoint obstruction in its
cleanest form.

\begin{lemma}[Boundary exponent collapse]
\label{lem:hardymeasure}
Let $0<r<s<\infty$ and let $\nu$ be a finite positive Borel measure on $\T$.
If
\[
 \|P\|_{L^s(\nu)}\le C\|P\|_{H^r}
 \qquad\hbox{for every analytic polynomial }P,
\]
then $\nu=0$.
\end{lemma}

\begin{proof}
Let $I\subset\T$ be an arc of normalized length $\delta\le1/2$, centered at
$\zeta$, and take $a=(1-\delta)\zeta$.  The analytic function
\[
 F_{a,r}(z)=\frac{(1-|a|^2)^{1/r}}{(1-\overline a z)^{2/r}}
\]
uses the analytic branch of the power in the disc.  Its boundary modulus to
the $r$th power is the Poisson kernel, hence $\|F_{a,r}\|_{H^r}=1$.
Moreover $|F_{a,r}|\ge c_r\delta^{-1/r}$ on $I$.  Since $F_{a,r}$ is
analytic beyond the closed disc, its Taylor polynomials converge uniformly
there; the assumed inequality therefore applies to $F_{a,r}$.  Consequently
\[
 \nu(I)^{1/s}c_r\delta^{-1/r}\le C,
 \qquad\text{so}\qquad
 \nu(I)\le C_{r,s}\delta^{s/r}.
\]
Partition $\T$ into $N$ equal arcs.  Since $s/r>1$,
\[
 \nu(\T)\le C_{r,s}N\,N^{-s/r}\longrightarrow0.
\]
Thus $\nu=0$.
\end{proof}

\begin{lemma}[Non-atomic $L^2\to L^p$ collapse]
\label{lem:nonatomic}
Let $p>2$, let $\mu$ be a finite non-atomic positive measure on $\T$, and
let $\nu$ be another finite positive measure.  If
\[
 \|b\|_{L^p(\nu)}\le C\|b\|_{L^2(\mu)}
 \tag{3.1}\label{eq:measureembed}
\]
for every Laurent polynomial $b$, then $\nu=0$.
\end{lemma}

\begin{proof}
Uniform density of Laurent polynomials in $C(\T)$ first extends
\eqref{eq:measureembed} to continuous functions.  This implies
$\nu\ll\mu$: if a compact $E$ has $\mu(E)=0$, choose continuous functions
equal to one on $E$ and supported in open sets of arbitrarily small
$\mu$-measure.  Regularity then handles all $\mu$-null sets.
Approximation in $L^2(\mu)+L^p(\nu)$ extends \eqref{eq:measureembed} to
indicators of Borel sets:
\[
 \nu(E)^{1/p}\le C\mu(E)^{1/2}.
 \tag{3.2}\label{eq:setembed}
\]
For each $N$, non-atomicity supplies a measurable partition
$\T=E_1\sqcup\cdots\sqcup E_N$ with
$\mu(E_j)=\mu(\T)/N$.  Summing \eqref{eq:setembed} to the $p$th power gives
\[
 \nu(\T)\le C^p\mu(\T)^{p/2}N^{1-p/2}\longrightarrow0.
\]
\end{proof}

\section{Bounded orbit vectors live on the contact set}

Put $h=1/g$.  The hypothesis and the open mapping theorem give
$|g(z)|>1$ in $\D$, so $h\in H^\infty$, $\|h\|_\infty\le1$, and
$|h(z)|<1$ in $\D$.  Let
\[
 A=\{\zeta\in\T:|h(\zeta)|=1\},
 \qquad R:H^q\longrightarrow L^q(A),\quad Rf=f|_A.
\]

\begin{lemma}[Contact factorization of bounded orbits]
\label{lem:factor}
If $x\in X$ and $\sup_{n\ge0}\|T^nx\|\le M$, then there is a unique
$\varphi\in L^p(A)$ such that $x=R'\varphi$, and
$\|\varphi\|_{L^p(A)}\le M$.
\end{lemma}

\begin{proof}
For $f\in H^q$ and $n\ge0$,
\[
 x(f)=(T^nx)(h^nf),
\]
because $g^nh^n=1$.  Therefore
\[
 |x(f)|\le M\|h^nf\|_{H^q}.
\]
Dominated convergence on the boundary gives
\[
 \lim_{n\to\infty}\|h^nf\|_{H^q}
 =\left(\int_A|f|^q\,dm\right)^{1/q}=\|Rf\|_{L^q(A)}.
 \tag{4.1}\label{eq:contactlimit}
\]
If $m(A)=0$, this already gives $x=0$.  Otherwise, $R$ is injective by
Hardy boundary uniqueness.  Its range is dense in $L^q(A)$.  Indeed, an
$L^p(A)$ annihilator, extended by zero to $\T$, has one-sided vanishing
Fourier coefficients and is therefore the boundary function of a Hardy
function (up to conjugation); it vanishes on $\T\setminus A$ and hence is
zero by the same uniqueness theorem.  Here $m(\T\setminus A)>0$: otherwise
$h$ would be an invertible inner function, hence a unimodular constant,
contrary to the nonconstancy of $g$.

It follows from \eqref{eq:contactlimit} that $x$ defines a bounded functional
of norm at most $M$ on the dense range of $R$.  Extend it to $L^q(A)$ and
use the linear bilinear-pairing realization $L^q(A)'=L^p(A)$.  This gives
$x=R'\varphi$ and the stated norm bound.
Density gives uniqueness.
\end{proof}

Using the bilinear pairing, direct calculation gives
\[
 T R'\varphi=R'(u\varphi),
 \qquad u=g|_A,\qquad |u|=1\ \text{a.e. on }A.
 \tag{4.2}\label{eq:contactintertwine}
\]
Here $R'\varphi(f)=\int_A f\varphi\,dm$.

\section{Proof of the main theorem}

For $\xi\in\mathcal H$,
\[
 \sup_{n\ge0}\|T^nJ\xi\|
 =\sup_{n\ge0}\|JU^n\xi\|
 \le\|J\|\,\|\xi\|.
\]
Lemma~\ref{lem:factor}, including uniqueness, therefore produces a bounded
linear map $K:\mathcal H\to L^p(A)$ such that
\[
 J=R'K,\qquad \|K\|\le\|J\|.
\]
Using $TJ=JU$, \eqref{eq:contactintertwine}, and injectivity of $R'$ gives
\[
 KU=M_uK.
 \tag{5.1}\label{eq:KU}
\]

Apply the Wold decomposition
$\mathcal H=\mathcal H_{\rm u}\oplus\mathcal H_{\rm s}$, where $U$ is
unitary on $\mathcal H_{\rm u}$ and a unilateral shift (of arbitrary
multiplicity) on $\mathcal H_{\rm s}$.

First let $\xi$ belong to the wandering subspace of the shift part and put
$\varphi=K\xi$.  The vectors $U^n\xi$ are orthogonal.  Hence, for every
analytic polynomial $P$,
\begin{align*}
 \|P(u)\varphi\|_{L^p(A)}
 &=\|KP(U)\xi\|_{L^p(A)} \\
 &\le\|K\|\,\|P(U)\xi\|_{\mathcal H}
 =\|K\|\,\|\xi\|\,\|P\|_{H^2}.
\end{align*}
Push the measure $|\varphi|^pdm|_A$ forward by $u:A\to\T$ and call the
result $\nu$.  The last display is exactly
$\|P\|_{L^p(\nu)}\le C\|P\|_{H^2}$.  Lemma~\ref{lem:hardymeasure}, with
$r=2$ and $s=p$, gives $\nu=0$, so $K\xi=0$.  Equation \eqref{eq:KU}
then makes $K$ vanish on all of $\mathcal H_{\rm s}$.

Now take $\xi\in\mathcal H_{\rm u}$ and put $\varphi=K\xi$.  Let $\mu_\xi$
be the spectral measure of $\xi$ for the unitary restriction of $U$,
and again let $\nu$ be the pushforward of $|\varphi|^pdm|_A$ by $u$.
Because $M_u$ and $U|_{\mathcal H_{\rm u}}$ are invertible,
\eqref{eq:KU} also intertwines their negative powers.  Thus, for every
Laurent polynomial $b$,
\[
 \|b\|_{L^p(\nu)}=\|b(u)\varphi\|_{L^p(A)}
 \le\|K\|\,\|b(U)\xi\|_{\mathcal H}
 =\|K\|\,\|b\|_{L^2(\mu_\xi)}.
\]
The measure $\mu_\xi$ is non-atomic: an atom would produce an eigenvector of
$U$, hence a one-dimensional invariant subspace.  Lemma~\ref{lem:nonatomic}
gives $\nu=0$, so $K\xi=0$.  Therefore $K=0$ on both Wold summands, and
$J=R'K=0$.  This proves Theorem~\ref{thm:main}.

For Corollary~\ref{cor:criterion}, the inclusion of $F$ into $X$ extends,
by the stronger-Hilbert-topology assumption, to a bounded nonzero map $J$
from the Hilbert completion of $F$ into $X$.  The extended isometry $U$
satisfies $TJ=JU$, contradicting Theorem~\ref{thm:main}.

\section{What this settles, and what it does not}

The theorem strictly strengthens the immediate observation that
$L^2(A)\not\hookrightarrow L^p(A)$ for $p>2$.  It rules out weighted or
singular contact measures, smooth or Sobolev regularizations, arbitrary
unitary spectral models, and also the pure unilateral-shift part of an
arbitrary isometric Hilbert model.  It therefore proves that the phrase
``fails miserably'' preceding Question 5.4 reflects a structural endpoint,
not a missing estimate in the displayed $L^2(A)$ construction.

It does \emph{not} prove that $T$ is not weakly hypercyclic.  Weak closure in
an infinite-dimensional Banach space is not sequential: a weakly dense orbit
may approach zero only through vectors of unbounded norm.  The contact
factorization controls bounded forward orbits and all isometric Hilbert cores,
but not such genuinely nonsequential behavior.  Eight focused upgrade
attempts tested smooth-contact series, finite $n$-weak reductions,
Hardy--Sobolev cores, direct interpolation from $H^2$, and reciprocal-growth
obstructions.  Each either falls under Theorem~\ref{thm:main} or stops at this
nonsequential weak-closure gap.  Question 5.4 remains open in this packet.

\section{Novelty and verification notes}

The run's registry, solution, attempt, and proof-gap indexes were searched for
the arXiv identifier, title, ``coanalytic Toeplitz'', ``weakly hypercyclic'',
and the $p>2$ formulation.  Bounded web searches through 17 August 2026 used
the exact Question 5.4 wording and close variants.  They found the source,
later work on norm/subspace hypercyclic Toeplitz operators, and a 2025 paper on
smooth-symbol norm hypercyclicity, but no stated resolution of Question 5.4
and no version of Theorem~\ref{thm:main}.  Novelty confidence is therefore
moderate, not definitive.

The proof was audited at the following pressure points: the duality orientation
in \eqref{eq:contactintertwine}; the dominated-convergence identity
\eqref{eq:contactlimit}; density of boundary restriction; extension of the
measure inequalities from polynomials to continuous/Borel tests; atomlessness
of unitary spectral measures; and both Wold components.  No computational
claim is used.

\paragraph{Human-review recommendation.}
Check especially the range-density sentence in Lemma~\ref{lem:factor}, the
bilinear duality convention leading to $u=g|_A$, and the passage from Laurent
polynomials to indicators in Lemma~\ref{lem:nonatomic}.  These are the only
places where a convention or a standard density fact is compressed.

\begin{thebibliography}{9}

\bibitem{Duren}
P.~Duren,
\emph{Theory of $H^p$ Spaces},
Academic Press, 1970.

\bibitem{Shkarin}
S.~Shkarin,
\emph{Orbits of coanalytic Toeplitz operators and weak hypercyclicity},
arXiv:1210.3191 (2012), especially Theorem~1.2, Theorem~1.9, and
Question~5.4.

\bibitem{Wold}
H.~Wold,
\emph{A Study in the Analysis of Stationary Time Series},
Almqvist \& Wiksell, 1938; see the standard Wold decomposition for Hilbert
space isometries.

\end{thebibliography}

\end{document}
